% q0023.tex — Robin Hood (Income Distribution / Poverty) [Numeric Answer]
\question{
  id        = {0023},
  title     = {Robin Hood},
  source    = {OBECON},
  year      = {2026},
  round     = {Seletiva IEO -- Phase A},
  language  = {English},
  topic     = {Micro},
  subtopic  = {Income Distribution / Poverty},
  type      = {numeric},
  statement = {%
    Consider the following income distribution for a small economy with 8
    individuals:

    \vspace{6pt}
    \begin{center}
    \begin{tabular}{lc}
      \toprule
      \textbf{Individual} & \textbf{Annual Income (in \$)} \\
      \midrule
      Alice   & 3,000  \\
      Bob     & 4,000  \\
      Carlos  & 5,000  \\
      Diana   & 9,000  \\
      Eva     & 21,000 \\
      Felipe  & 25,000 \\
      Gabriel & 35,000 \\
      Helena  & 64,000 \\
      \bottomrule
    \end{tabular}
    \end{center}
    \vspace{6pt}

    The government implements a redistributive policy where the top 25\% of
    earners are taxed \$10,000 which are distributed equally among the bottom
    50\% of earners. If the poverty line is defined as 50\% of the
    \emph{status quo} median income, by what percentage does this policy
    reduce the number of people living in poverty?

    \textit{Note:} We define the median of an ordered sequence
    $a_1, a_2, \ldots, a_{2n}$ as $\dfrac{a_n + a_{n+1}}{2}$.

    Give the answer as a number rounded to the nearest integer. If the
    fractional part is 0.5, round to the nearest even integer.
  },
  choices   = {},
  answer    = {},
  solution  = {%
    % TODO: add solution
  },
}
